% =====================================================
% 题型：双曲线离心率填空题 (q_hyperbola_eccentricity.tex)
% 知识板块：双曲线
% 主思维：定义优先思维
% 副思维：公式结构
% 错因标签：标准式化简错误、a, b, c 方程记混
% 讲义用途：圆锥曲线基础
% =====================================================
% 难度：★ (基础)
% =====================================================
\newcommand{\qHyperbolaEccentricityStem}{双曲线 \(5x^2 - 6y^2 = 1\) 的离心率为 \blank{3cm}。}

\newcommand{\qHyperbolaEccentricityAnswer}{\(\sqrt{\dfrac{11}{6}}\)}

\newcommand{\qHyperbolaEccentricitySolution}{%
  \textbf{【解题思路】}\par
  求双曲线的离心率，核心在于将一般式方程化为标准形式 \(\dfrac{x^2}{a^2} - \dfrac{y^2}{b^2} = 1\)，进而读取参数 \(a^2, b^2\)，结合恒等式 \(c^2 = a^2 + b^2\) 求得焦距 \(c\)，最终由离心率定义 \(e = \dfrac{c}{a}\) 求解。\par\vspace{0.4em}

  \textbf{【解析计算】}\par
  已知双曲线方程为 \(5x^2 - 6y^2 = 1\)，将其化为标准形式：
  \[
    \frac{x^2}{\frac{1}{5}} - \frac{y^2}{\frac{1}{6}} = 1.
  \]
  由此可得焦点在 \(x\) 轴上，且参数为：
  \[
    a^2 = \frac{1}{5}, \qquad b^2 = \frac{1}{6}.
  \]
  在双曲线中，焦距平方满足关系式：
  \[
    c^2 = a^2 + b^2 = \frac{1}{5} + \frac{1}{6} = \frac{11}{30}.
  \]
  由离心率定义公式，可得：
  \[
    \begin{aligned}
      e^2 &= \frac{c^2}{a^2} = \frac{\frac{11}{30}}{\frac{1}{5}} = \frac{11}{6} \\[0.4em]
      \implies e &= \sqrt{\frac{11}{6}} \quad \text{（或化简为 } \frac{\sqrt{66}}{6}\text{）}.
    \end{aligned}
  \]
  故答案为 \(\sqrt{\dfrac{11}{6}}\)。%
}

\newcommand{\qHyperbolaEccentricityDiagnosis}{%
  学生在化简方程时，容易将 \(5x^2\) 变形错误（如错记为 \(a^2 = 5\)）。同时，双曲线的几何性质中 \(c^2 = a^2 + b^2\) 极易与椭圆的 \(c^2 = a^2 - b^2\) 记混。%
}

\newcommand{\qHyperbolaEccentricityTeachingNotes}{%
  离心率计算是圆锥曲线的基础考点。讲评时应重点强化两点：一是“化标准”，将二次项系数化为 1 后，分母即为对应的平方项；二是区分椭圆与双曲线的 \(a, b, c\) 方程关系，夯实基础。%
}
